% Web専用初稿．独立査読・文献該当箇所照合未了．
\section{非Archimedes絶対値}
% genron:begin K-001661
% genron:use K-000001
% genron:use K-000002
% genron:use K-000003
% genron:use K-000004
% genron:use K-000005
% genron:use K-000007
% genron:use K-000017
% genron:use K-000024
% genron:use K-001591
% genron:use K-001594
\begin{defn}[非Archimedes絶対値]
体 $K$ の絶対値 $|\cdot|$ が $|x+y|\leq\max(|x|,|y|)$ を満たすとき非Archimedes的という．
\end{defn}
% genron:end K-001661
% genron:begin K-001662
% genron:use K-000001
% genron:use K-000002
% genron:use K-000003
% genron:use K-000004
% genron:use K-000005
% genron:use K-000007
% genron:use K-000017
% genron:use K-000024
% genron:use K-001661
\begin{prop}[整数倍による判定]
絶対値が非Archimedes的であることと，全ての正整数 $n$ について $|n1_K|\leq1$ であることは同値である．
\end{prop}
\begin{proof}
必要性は強三角不等式による帰納法で従う．逆に二項展開により $|x+y|^m\leq(m+1)\max(|x|,|y|)^m$ を得る．各二項係数の絶対値は仮定により高々 $1$ である．$m$ 乗根を取り $m\to\infty$ とすれば強三角不等式を得る．
\end{proof}
% genron:end K-001662
% genron:begin K-001663
% genron:use K-000001
% genron:use K-000002
% genron:use K-000003
% genron:use K-000004
% genron:use K-000005
% genron:use K-000007
% genron:use K-000017
% genron:use K-000024
% genron:use K-001661
\begin{prop}[異なる絶対値を持つ元の和]
$|x|<|y|$ なら $|x+y|=|y|$ である．
\end{prop}
\begin{proof}
強三角不等式から $|x+y|\leq |y|$ である．一方 $y=(x+y)-x$ より $|y|\leq\max(|x+y|,|x|)$ である．$|x|<|y|$ なので逆の不等式を得る．
\end{proof}
% genron:end K-001663
% genron:begin K-001664
% genron:use K-000001
% genron:use K-000002
% genron:use K-000003
% genron:use K-000004
% genron:use K-000005
% genron:use K-000007
% genron:use K-000017
% genron:use K-000024
% genron:use K-001661
\begin{defn}[付値環とその極大イデアル]
非Archimedes絶対値付き体 $K$ に対し $\mathcal O=\{x:|x|\leq1\}$，$\mathfrak m=\{x:|x|<1\}$ と置く．
\end{defn}
% genron:end K-001664
% genron:begin K-001665
% genron:use K-000001
% genron:use K-000002
% genron:use K-000003
% genron:use K-000004
% genron:use K-000005
% genron:use K-000007
% genron:use K-000017
% genron:use K-000024
% genron:use K-001664
% genron:use K-000446
\begin{prop}[付値環の局所性]
$\mathcal O$ は局所整域であり，唯一の極大イデアルは $\mathfrak m$ である．任意の $x\in K^\times$ に対し $x$ 又は $x^{-1}$ が $\mathcal O$ に属する．
\end{prop}
\begin{proof}
強三角不等式と乗法性により $\mathcal O$ は部分環，$\mathfrak m$ は真のイデアルである．$\mathcal O$ の元が単元であることはその絶対値が $1$ であることと同値である．従って任意の真のイデアルは $\mathfrak m$ に含まれる．最後の主張は $|x^{-1}|=|x|^{-1}$ から従う．
\end{proof}
% genron:end K-001665
% genron:begin K-001666
% genron:use K-000001
% genron:use K-000002
% genron:use K-000003
% genron:use K-000004
% genron:use K-000005
% genron:use K-000007
% genron:use K-000017
% genron:use K-000024
% genron:use K-001665
\begin{defn}[剰余体と値群]
$\mathcal O/\mathfrak m$ を剰余体といい，乗法群 $|K^\times|\subseteq\mathbb R_{>0}$ を値群という．
\end{defn}
% genron:end K-001666
% genron:begin K-001667
% genron:use K-000001
% genron:use K-000002
% genron:use K-000003
% genron:use K-000004
% genron:use K-000005
% genron:use K-000007
% genron:use K-000017
% genron:use K-000024
% genron:use K-001665
\begin{prop}[付値環の主イデアルの比較]
$a,b\in\mathcal O$ が非零なら $a\mathcal O\subseteq b\mathcal O$ と $|a|\leq |b|$ は同値である．従って任意の二イデアルは包含関係で比較可能である．
\end{prop}
\begin{proof}
最初の同値は $a/b\in\mathcal O$ と同値であることから従う．二イデアル $I,J$ が比較不能なら $a\in I\setminus J$ と $b\in J\setminus I$ が取れる．主イデアルの比較可能性は $a\in J$ 又は $b\in I$ を強制し矛盾する．
\end{proof}
% genron:end K-001667
